\chapter{Virtual memory, demand paging, and interrupts}
\label{cap:memoria-usuario}

This chapter documents the \code{gdt.rs}, \code{idt.rs},
\code{paging.rs}, and \code{user.rs} modules, added after the console and the shell.
The goal is to move away from the ``ring 0, interrupts disabled, and
static memory'' model to a kernel capable of:

\begin{itemize}
  \item receiving exceptions and IRQs through an IDT;
  \item receiving the PIT timer on IRQ0;
  \item allocating 4 KiB physical frames;
  \item materializing pages only when the CPU generates \code{\#PF};
  \item executing code in ring 3 with a real barrier between user and
        kernel.
\end{itemize}

\section{Initialization order}

The sequence in \code{kernel\_main} became:

\begin{enumerate}
  \item \code{vga::init()} and \code{serial::init()};
  \item \code{paging::init(magic, info)} --- frees the available frames
        described by Multiboot2;
  \item \code{gdt::init()} --- final GDT, ring 3 segments, and TSS;
  \item \code{keyboard::init()} --- PS/2 keyboard by \emph{polling};
  \item \code{idt::init()} --- IDT, PIC remapping, and PIT;
  \item \code{sti} and entry into the shell.
\end{enumerate}

The IDT is installed only after the GDT has the TSS, because the
\emph{double fault} gate uses an IST stack. The timer is unmasked only after
the IDT is ready.

\section{GDT, TSS, and privilege stacks}

The GDT of the boot stub contains only ring 0 descriptors. To enter
ring 3, \code{gdt.rs} builds a table with seven entries:

\begin{table}[H]
\centering
\caption{Descriptors of the final GDT.}
\label{tab:gdt-usuario}
\begin{tabular}{@{}rll@{}}
\toprule
\textbf{Selector} & \textbf{Contents} & \textbf{DPL} \\
\midrule
\code{0x08} & 64-bit code & 0 \\
\code{0x10} & data & 0 \\
\code{0x18} & data & 3 \\
\code{0x20} & 64-bit code & 3 \\
\code{0x28} & TSS (two entries) & 0 \\
\bottomrule
\end{tabular}
\end{table}

The TSS is mandatory for automatic stack switching. When the CPU
moves from ring 3 to ring 0 due to an interrupt or exception, it reads
\code{rsp0} from the TSS and continues pushing onto that stack. The kernel reserves:

\begin{itemize}
  \item \code{USER\_KERNEL\_STACK}: 16 KiB for interrupts coming from
        user mode;
  \item \code{DOUBLE\_FAULT\_STACK}: 8 KiB as \code{IST1}, so that a
        \code{\#DF} does not try to use a possibly corrupted stack.
\end{itemize}

The \code{iomap\_base} field points to the end of the TSS, disabling the I/O
bitmap. Thus, \code{in}/\code{out} instructions in ring 3 generate \code{\#GP}.

\section{IDT, stubs, and error-code normalization}

The IDT has 256 gates of 16 bytes. The first 32 are exceptions; the
next 16 are the PIC IRQs remapped to \code{0x20..0x2F}; gate
\code{0x80} is a didactic \emph{syscall} with DPL = 3.

The classic problem is that the CPU pushes an error code only on
some exceptions (8, 10--14, 17, 21, 29, and 30). If each \emph{stub}
handled the stack differently, the layout seen by Rust would change. To
standardize, \code{idt.rs} generates stubs in assembly:

\begin{lstlisting}[caption={Stubs that normalize the error code.},label={lst:isrstubs}]
.macro isr_no_err n
.global isr_stub_\n
isr_stub_\n:
    pushq $0          # erro sintetico
    pushq $\n         # numero do vetor
    jmp isr_common
.endm

.macro isr_err n
.global isr_stub_\n
isr_stub_\n:
    pushq $\n         # erro real ja empilhado pela CPU
    jmp isr_common
.endm
\end{lstlisting}

The common \emph{stub} saves the general-purpose registers, preserves
\code{xmm0..xmm15} and \code{MXCSR}, aligns the stack to call
\code{interrupt\_dispatch}, and, on return, restores everything and executes
\code{iretq}. The frame seen by Rust is:

\begin{lstlisting}[caption={Layout of \code{InterruptFrame} after the common stub.},label={lst:frame}]
#[repr(C)]
pub struct InterruptFrame {
    pub r15: u64, /* ... r8 */
    pub rdi: u64, /* ... rax */
    pub vector: u64,
    pub error: u64,
    pub rip: u64,
    pub cs: u64,
    pub rflags: u64,
    pub rsp: u64,   // so existe se veio de ring 3
    pub ss: u64,    // so existe se veio de ring 3
}
\end{lstlisting}

Vector 14 calls the \emph{page fault} handler; vector 13 handles
\code{\#GP}; vector 0x80 terminates the user program; vector 32
increments the \emph{ticks} counter and sends EOI to the PIC.

\section{PIC and PIT}

The 8259 PIC is remapped so as not to collide with the exceptions:

\begin{lstlisting}[caption={PIC remapping and PIT initialization.},label={lst:picpit}]
outb(0x20, 0x11); outb(0xA0, 0x11);  // ICW1
outb(0x21, 0x20); outb(0xA1, 0x28);  // ICW2: vetores 0x20/0x28
outb(0x21, 0x04); outb(0xA1, 0x02);  // ICW3
outb(0x21, 0x01); outb(0xA1, 0x01);  // ICW4

let divisor = 11_931;                // 1,193182 MHz / 100 Hz
outb(0x43, 0x36);
outb(0x40, divisor as u8);
outb(0x40, (divisor >> 8) as u8);

outb(0x21, 0xFE);                    // so IRQ0
outb(0xA1, 0xFF);
\end{lstlisting}

IRQ1 remains masked because \code{keyboard.rs} reads the 8042 by
\emph{polling}. Enabling it would require a scancode buffer and coordinating
access to the controller.

\section{Physical frame allocator}

\code{paging::init} walks the Multiboot2 tags. Tag 6 (memory map)
provides 24-byte entries with base, size, and type. The allocator starts
with all 262,144 frames of the first GiB marked as used and frees
only the ranges of type 1 (available RAM).

The bitmap is 32 KiB:

\begin{lstlisting}[caption={Frame bitmap and allocation.},label={lst:bitmap}]
static mut FRAME_BITMAP: [u8; 32 * 1024] = [0xFF; 32 * 1024];

pub fn alloc_frame() -> Option<u64> {
    /* procura um bit zero, marca como usado e zera os 4096 bytes */
}
\end{lstlisting}

The first MiB and the end of the kernel (symbol \code{\_\_kernel\_end} from the
\emph{linker script}) are always reserved. The Multiboot2 structure is also
reserved after the available ranges are freed.

Since the stub maps the first GiB by identity, writing to the physical
address of the frame is equivalent to writing to the low virtual address.
This makes it possible to zero frames without needing an additional direct map.

\section{Demand paging and U/S protection}

Each table entry has present (P), writable (W), and user (U/S) bits. x86\_64
requires U/S = 1 at \emph{all} levels (PML4, PDPT, PD, and PT) to allow an
access in ring 3. The kernel uses this to create two regions:

\begin{table}[H]
\centering
\caption{Demand paging regions.}
\label{tab:demand}
\begin{tabularx}{\linewidth}{@{}lrrX@{}}
\toprule
\textbf{Region} & \textbf{Base} & \textbf{Size} & \textbf{Leaf U/S} \\
\midrule
User & \code{0x0000008000000000} & 2 MiB & 1 \\
Kernel   & \code{0xFFFFFFFFC0000000} & 2 MiB & 0 \\
\bottomrule
\end{tabularx}
\end{table}

No leaf of these regions exists at boot. When \code{\#PF} occurs,
the handler reads \code{CR2} and the error code and decides:

\begin{enumerate}
  \item if the address belongs to the user region and the page is not
        present, it allocates a frame and maps it with
        \code{P|W|U/S}; the intermediate tables also receive U/S = 1;
  \item if it belongs to the kernel region and the error's U/S bit is zero,
        it maps with \code{P|W} and U/S = 0;
  \item if a user tries the kernel region, the error has U/S = 1 and the
        handler refuses the mapping; the user program is terminated;
  \item if the address is outside both regions, the fault is fatal
        to the kernel or terminates the user.
\end{enumerate}

\begin{nota}[Why the kernel cannot ``open up'' a supervisor table]
Mapping a user page is not just a matter of setting U/S on the leaf.
If any intermediate level is supervisor, the hardware denies the access.
Moreover, turning a kernel table into a user table could expose other mappings.
For this reason \code{next\_table} refuses a supervisor table on the path of a
user page.
\end{nota}

\section{User mode and return to kernel}

\code{user.rs} builds a small machine program in a buffer and copies it
to the \code{USER\_CODE} page, which is resident. The stack and the data
page are on demand:

\begin{enumerate}
  \item the program performs \code{push}/\code{pop}, generating \code{\#PF} on the stack;
  \item it writes to \code{USER\_DATA}, generating another \code{\#PF};
  \item it executes \code{int \$0x80} with \code{rax = 0}.
\end{enumerate}

The \code{userprot} command uses another program: it tries to read
\code{0xFFFFFFFF80000000} (start of the kernel). The page exists, but it has
U/S = 0; the \emph{page fault} error has P = 1 and U/S = 1. The handler
recognizes the violation, prints the data, and does not return control to the
program.

Entry into ring 3 is done by \code{iretq}, which loads CS/SS from DPL 3,
RIP, RSP, and RFLAGS. The return does not use the normal \code{iretq} path:

\begin{itemize}
  \item \code{user\_enter} saves the callee-saved registers and the
        kernel RSP in \code{user\_saved\_rsp};
  \item an exception or \code{int \$0x80} lands in the common \emph{stub} in
        ring 0;
  \item \code{user::abort} writes the status and calls \code{user\_resume}, which
        restores the saved RSP, reloads the kernel segments, and returns to
        the \code{user\_enter} call.
\end{itemize}

This mechanism is a controlled \emph{longjmp}. It avoids having to reconstruct
the interrupt frame to return to ring 0 and keeps the shell simple.

\section{Commands and evidence}

\begin{table}[H]
\centering
\caption{Commands that exercise the new infrastructure.}
\label{tab:comandos-memoria}
\begin{tabularx}{\linewidth}{@{}lX@{}}
\toprule
\textbf{Command} & \textbf{Effect} \\
\midrule
\code{demand} & writes to the kernel demand paging region; the first
access generates \code{\#PF} \\
\code{usertest} & enters ring 3, generates \code{\#PF} on the stack and on the data, and
exits via \code{int \$0x80} \\
\code{userprot} & enters ring 3 and tries to read the kernel; U/S protection
terminates the program \\
\code{mem} & shows allocated/free frames, \emph{page faults}, \emph{ticks},
and the U/S bits of a user page and of the kernel \\
\bottomrule
\end{tabularx}
\end{table}

A typical execution in QEMU:

\begin{lstlisting}[style=shell,caption={Output of the memory and ring 3 commands.},label={lst:testdemand}]
rinux> demand
[PF] demanda atendida: 0xffffffffc0000000 (kernel, erro=0x2)
[DEMAND] 0xffffffffc0000000 = 0xdeadbeefcafe1234 (page faults: 1)
rinux> usertest
[USER] entrando em ring 3 (demand paging)...
[PF] demanda atendida: 0x0000008000007ff8 (usuario, erro=0x6)
[PF] demanda atendida: 0x0000008000002000 (usuario, erro=0x6)
[USER] voltou ao kernel; page faults nesta execucao: 2
rinux> userprot
[USER] entrando em ring 3 (teste de protecao)...
[USER] page fault em 0xffffffff80000000: endereco fora das regioes de demand paging (erro=0x5, rip=0x000000800000100a)
[USER] programa encerrado com status 1
rinux> mem
Frames fisicos (4 KiB):
  alocados : 7
  livres   : 32217
  page faults atendidos: 3
  ticks do timer       : 372
  USER_CODE  flags=0x203067 (U/S=1)
  KERNEL_VMA flags=0xe3 (U/S=0)
\end{lstlisting}

The \code{0x6} error (P = 0, W = 1, U/S = 1) confirms a user write to an
absent page; the \code{0x5} error (P = 1, U/S = 1, without W) confirms a
user read from a present, supervisor page.

\section{Limitations of this implementation}

\begin{itemize}
  \item there are no separate address spaces per process: the
        user program uses the same PML4 as the kernel;
  \item preemption is disabled while a ring 3 program is active, because it
        uses the shared TSS interrupt stack; there is no per-process kernel
        stack or \code{fork};
  \item the syscalls are the original four plus the file and PID calls;
        user pointers are range-checked, but there is no
        \code{copy\_from\_user} helper;
  \item SMEP/SMAP and NX are not yet enabled; the separation
        depends only on U/S;
  \item the keyboard remains in \emph{polling} and IRQ1 remains masked;
  \item the demand paging regions have a fixed size of 2 MiB and lie
        outside the first GiB mapped by identity.
\end{itemize}
